\documentclass[a4paper,10pt]{article}

\usepackage[english]{babel}
\usepackage{graphicx}
\usepackage[colorlinks, linkcolor=black, citecolor=black, urlcolor=black]{hyperref}
\usepackage{geometry}
\geometry{tmargin=3cm, bmargin=2.2cm, lmargin=2.2cm, rmargin=2cm}
\usepackage{todonotes} %Used for the figure placeholders

% Your name and student number must be filled in on the title page found in
% titlepage.tex.

\begin{document}
\begin{titlepage}
    \newpage
    \thispagestyle{empty}
    \frenchspacing
    \hspace{-0.2cm}
    \includegraphics[height=3.4cm]{sedes}
    \hspace{0.2cm}
    \rule{0.5pt}{3.4cm}
    \hspace{0.2cm}
    \begin{minipage}[b]{8cm}
        \Large{Katholieke\newline Universiteit\newline Leuven}\smallskip\newline
        \large{}\smallskip\newline
        \textbf{Department of\newline Computer Science}\smallskip
    \end{minipage}
    \hspace{\stretch{1}}
    \vspace*{3.2cm}\vfill
    \begin{center}
        \begin{minipage}[t]{\textwidth}
            \begin{center}
                \LARGE{\rm{\textbf{\uppercase{Document Processing}}\\ADD
                application}}\\
                \Large{\rm{Software Architecture (H09B5a and H07Z9a) -- 
                Part 2a}}
            \end{center}
        \end{minipage}
    \end{center}
    \vfill
    \hfill\makebox[8.5cm][l]{%
        \vbox to 7cm{\vfill\noindent
            {\rm \textbf{Student A (r123456)}}\\
            {\rm \textbf{Student B (r987654)}}\\[2mm]
            {\rm Academic year 2014--2015}
        }
    }
\end{titlepage}


\tableofcontents
\newpage

\section{Domain analysis}\label{sec:domain}
\subsection{Domain models}
This section shows the domain model(s).

\begin{figure}[!htp]
    \centering
    %\includegraphics[width=0.8\textwidth]{}
    \missingfigure[figwidth=0.8\textwidth]{Domain model}
    \caption{The domain model for the system.}\label{fig:domain_model}
\end{figure}

\subsection{Domain constraints}
In this section we provide additional domain constraints.

\begin{itemize}
    \item This is a first constraint.
    \item This is a second constraint.
\end{itemize}

\subsection{Glossary}
In this section, we provide a glossary of the most important terminology used
in this analysis.

\begin{itemize}
    \item \textbf{Term1}: definition
    \item \textbf{Term2}: definition
\end{itemize}

\section{Functional requirements}\label{sec:functional}
\subsection*{Use case model}

\begin{figure}[!htp]
    \centering
    %\includegraphics[width=0.8\textwidth]{}
    \missingfigure[figwidth=0.8\textwidth]{Use case model}
    \caption{Use case diagram for the system.}\label{fig:use_case_model}
\end{figure}

\subsection{\emph{UC1}: Name of use case 1}
\begin{itemize}
    \item \textbf{Name:} Name of use case 1
    \item \textbf{Primary actor:} primary actor
    \item \textbf{Interested parties:} 
        \begin{itemize}
            \item \textit{Name of interested party:} reason why party is interested
        \end{itemize}

    \item \textbf{Preconditions:}
        \begin{itemize}
            \item First precondition.
            \item Second precondition.
        \end{itemize}

    \item \textbf{Postconditions:}
        \begin{itemize}
            \item First postcondition.
            \item Second postcondition.
        \end{itemize}
        
    \item \textbf{Main scenario:} 
    \begin{enumerate}
       \item Step 1
       \item Step 2
       \item Step 3
       \item \ldots
    \end{enumerate}

    \item \textbf{Alternative scenarios:} 
    \begin{enumerate}
        \item [3b.] Alternative at step 3
    \end{enumerate}
    
    \item \textbf{Remarks:}
        \begin{itemize}
            \item First remark
        \end{itemize}
\end{itemize}

\subsection{\emph{UC1.1}: Consult document status}
\begin{itemize}
    \item \textbf{Name:} Consult document status
    \item \textbf{Primary actor:} Customer organisation
    \item \textbf{Interested parties:} 
        \begin{itemize}
            \item \textit{Customer organisation:} wants to consult the status of one of its documents at eDocs
            \item \textit{The eDocs system:} provides a management dashboard through which the customer can, among other things, consult the status of its documents
        \end{itemize}

    \item \textbf{Preconditions:}
        \begin{itemize}
            \item The customer organisation is a registred client at eDocs
            \item The customer organisation 
            \item Second precondition.
        \end{itemize}

    \item \textbf{Postconditions:}
        \begin{itemize}
            \item First postcondition.
            \item Second postcondition.
        \end{itemize}
        
    \item \textbf{Main scenario:} 
    \begin{enumerate}
       \item Step 1
       \item Step 2
       \item Step 3
       \item \ldots
    \end{enumerate}

    \item \textbf{Alternative scenarios:} 
    \begin{enumerate}
        \item [3b.] Alternative at step 3
    \end{enumerate}
    
    \item \textbf{Remarks:}
        \begin{itemize}
            \item First remark
        \end{itemize}
\end{itemize}

\subsection{\emph{UC2}: Register customer organisation}
\begin{itemize}
    \item \textbf{Name:} Register customer organisation
    \item \textbf{Summary:} The eDocs administrator registers a new customer organisation in order to enlarge the eDocs client base
    \item \textbf{Primary actor:} eDocs administrator
    \item \textbf{Secondary actor:} Customer organisation
    \item \textbf{Interested parties:} 
        \begin{itemize}
            \item \textit{Customer organisation:} wants to outsource its document processing to eDocs
            \item \textit{eDocs:} wants to gain new clients  customer organisations) for which it can process documents
            \item \textit{eDocs administrator:} deals with the administrational tasks during the registration of customer organisations
        \end{itemize}
    \item \textbf{Preconditions:}
        \begin{itemize}
            \item The customer organisation has actively expressed its interest in eDocs
            \item The eDocs administrator has initiated negotiations about the contract and SLA with the customer organisation
        \end{itemize}
    \item \textbf{Postconditions:}
        \begin{itemize}
            \item eDocs has gained a new client (and friend)
            \item The customer organisation and eDocs administrator have succesfully negotiated a contract and SLA
        \end{itemize}
        
    \item \textbf{Main scenario:} 
    \begin{enumerate}
       \item The customer organisation has expressed its interest in eDocs as its potential document processing partner
       \item The eDocs administrator starts negotiations with the customer organisation about the contract and SLA
       \item Both parties agree upon the terms and conditions specified in the forenamed documents
       \item The customer organisation is now a client of eDocs and has agreed to outsource (part of) its document processing tasks to eDocs
       \item \ldots
    \end{enumerate}
    \item \textbf{Alternative scenarios:} 
    \begin{enumerate}
        \item [3a.] The two parties do not reach an agreement
        \item [4a.] The customer organisation expresses its intent to outsource its document processing tasks elsewhere
        \item [5a.] eDocs had lost a potential client
    \end{enumerate}
    
    \item \textbf{Assumptions:}
        \begin{itemize}
            \item The registration process is initiated by the customer organisation expressing its interest in eDocs. It is thus assumed that eDocs is a big player in the document processing market, following a passive strategy to gain new clients.
        \end{itemize}
\end{itemize}


\subsection{\emph{UC3}: Consult management dashboard}
\begin{itemize}
    \item \textbf{Name:} Consult management dashboard
    \item \textbf{Summary:} 
    \item \textbf{Primary actor:} Customer organisation
    \item \textbf{Secondary actor:} eDocs system
    \item \textbf{Interested parties:} 
        \begin{itemize}
            \item \textit{Customer organisation:} wants to access and/or change some of the organisation's data that is stored in the eDocs system
            \item \textit{System:} provides an interface to access and manipulate the organisation's data that is outsourced to eDocs in a consistent way
        \end{itemize}
    \item \textbf{Preconditions:}
        \begin{itemize}
            \item The customer organisation is a registered client of eDocs
            \item The customer organisation is logged in to the eDocs system
        \end{itemize}
    \item \textbf{Postconditions:}
        \begin{itemize}
            \item The customer organisation has accessed and/or manipulated (part of) its data that is stored/processed by eDocs
            \item 
        \end{itemize}
        
    \item \textbf{Main scenario:} 
    \begin{enumerate}
       \item The customer organisation has expressed its interest in eDocs as its potential document processing partner
       \item The eDocs administrator starts negotiations with the customer organisation about the contract and SLA
       \item Both parties agree upon the terms and conditions specified in the forenamed documents
       \item The customer organisation is now a client of eDocs and has agreed to outsource (part of) its document processing tasks to eDocs
       \item \ldots
    \end{enumerate}
    \item \textbf{Alternative scenarios:} 
    \begin{enumerate}
        \item [3a.] The two parties do not reach an agreement
        \item [4a.] The customer organisation expresses its intent to outsource its document processing tasks elsewhere
        \item [5a.] eDocs had lost a potential client
    \end{enumerate}
    
    \item \textbf{Remarks:}
        \begin{itemize}
            \item 
        \end{itemize}
\end{itemize}


\subsection{\emph{UC4}: Deliver raw data}
\begin{itemize}
    \item \textbf{Name:} Deliver raw data
    \item \textbf{Summary:} The customer organisation provides eDocs with the raw data to be processed
    \item \textbf{Primary actor:} Customer organisation
    \item \textbf{Secondary actor:} eDocs system
    \item \textbf{Interested parties:} 
        \begin{itemize}
            \item \textit{Customer organisation:} wants to deliver its raw data to eDocs so that the latter can process it accordingly
            \item \textit{System:} accepts the raw data to be processed
        \end{itemize}
    \item \textbf{Preconditions:}
        \begin{itemize}
            \item The customer organisation is logged in to the eDocs system
        \end{itemize}
    \item \textbf{Postconditions:}
        \begin{itemize}
            \item The customer organisation has provided the eDocs system with its raw data
            \item The eDocs system has succesfully received the customer organisation's raw data
        \end{itemize}
        
    \item \textbf{Main scenario:} 
    \begin{enumerate}
       \item The customer organisation sends its raw data to eDocs via the management dashboard or via a web service that provides a link from the customer organisation's system to that of eDocs
       \item The eDocs system receives the customer organisation's raw data
       \item \ldots
    \end{enumerate}
\end{itemize}


\subsection{\emph{UC5}: Process raw data}
\begin{itemize}
    \item \textbf{Name:} Process raw data
    \item \textbf{Summary:} The eDocs system process the raw data that was provided by the customer organisation in a previous step
    \item \textbf{Primary actor:} eDocs system
    \item \textbf{Secondary actor:} None
    \item \textbf{Interested parties:} 
        \begin{itemize}
            \item \textit{Customer organisation:} wants its raw data to be processed correctly
            \item \textit{eDocs system:} wants to process the raw data correctly
        \end{itemize}
    \item \textbf{Preconditions:}
        \begin{itemize}
            \item The customer organisation, to which the raw data belongs, is a registered client of eDocs
            \item The raw data that has to be processed, was succesfully delivered to the eDocs system
        \end{itemize}
    \item \textbf{Postconditions:}
        \begin{itemize}
            \item The raw data is processed
        \end{itemize}
        
    \item \textbf{Main scenario:} 
    \begin{enumerate}
       \item Check raw data for consistency and completeness
       \item Process document data according to previously provided template
       \item Process meta data
       \item \ldots
    \end{enumerate}
    \item \textbf{Alternative scenarios:} 
    \begin{enumerate}
        \item [2a.] Erronous data is sent back to the customer organisation from which it originated
        \item [3a.] The customer organisation corrects the erronous part(s) of the respective (batch of) raw data
        \item [4a.] The customer organisation delivers a (batch of) raw data consisting only of those corrected records that were (partly) erronous at the beginning of this processing phase
        \item [5a.] Restart from step 1.
    \end{enumerate}


\subsection{\emph{UC6}: De-activate personal document store}
\begin{itemize}
	\item \textbf{Summary:} A Registered recipient de-activates his or her personal document store account.
    \item \textbf{Primary actor:} Registered recipient
    \item \textbf{Secondary actors:} The System
    \item \textbf{Interested parties:}
        \begin{itemize}
            \item \textit{Registered recipient:} wants to de-activite his or her personal document store
            \item \textit{eDocs:} wants to keep as many Registered recipients as possible.
        \end{itemize}

    \item \textbf{Preconditions:} None

    \item \textbf{Postconditions:}
        \begin{itemize}
            \item The Registered recipient is now an Unregistered recipient and can no longer log in to the system.
            \item The Registered recipient is now logged out and can no longer make use of his/her personal document store.
            \item The System has removed the account of the Registered Recipient. It keeps the documents of the Registered recipient archived.
            \item The Registered recipient now receives his or her documents via another channel than his or her personal document store.
        \end{itemize}
        
    \item \textbf{Main scenario:} 
    \begin{enumerate}
       \item The Registered recipient indicates that he or she wants to de-activate his or her personal document store.
       \item The System removes the account of the registered recipient.
       \item The System indicates success to the Registered recipient and logs him or her out.
    \end{enumerate}

    \item \textbf{Alternative scenarios:} 
    \begin{enumerate}
        \item \ldots
    \end{enumerate}
    
    \item \textbf{Remarks:}
        \begin{itemize}
        	\item None
        \end{itemize}
\end{itemize}


\subsection{\emph{UC7}: Provide document template}
\begin{itemize}
	\item \textbf{Summary:} An eDocs administrator
    \item \textbf{Primary actor:} eDocs administrator
    \item \textbf{Interested parties:} 
        \begin{itemize}
        	\item \textit{eDocs administrator:} wants to provide a document template that a customer organization can fill in.
            \item \textit{Customer organization:} wants to have a template for its documents
            \item \textit{eDocs:} wants to generate documents based on a template that the Customer organization has filled in.
        \end{itemize}

    \item \textbf{Preconditions:}
        \begin{itemize}
            \item The Customer organization has registered itself with eDocs.
        \end{itemize}

    \item \textbf{Postconditions:}
        \begin{itemize}
            \item The System now contains a template which the Customer organization can fill in.
        \end{itemize}
        
    \item \textbf{Main scenario:} 
    \begin{enumerate}
		\item The eDocs administrator provides a document template to the System.
    \end{enumerate}

    \item \textbf{Alternative scenarios:} 
    \begin{enumerate}
    	\item None
    \end{enumerate}
    
    \item \textbf{Remarks:}
        \begin{itemize}
        	\item None
        \end{itemize}
\end{itemize}

\subsection{\emph{UC8}: Update document template}
\begin{itemize}
	\item \textbf{Summary:} The Customer organization updates a document template.
    \item \textbf{Primary actor:} Customer organization
    \item \textbf{Secondary actors:} eDocs administrator
    \item \textbf{Interested parties:} 
        \begin{itemize}
            \item \textit{Customer organization:} wants its documents to have a layout corresponding to the initial template it fills in.
            \item \textit{eDocs:} wants to generate documents with a layout corresponding to the wishes of its clients.
        \end{itemize}

    \item \textbf{Preconditions:}
        \begin{itemize}
            \item The Customer organization is logged in.
        \end{itemize}

    \item \textbf{Postconditions:}
        \begin{itemize}
            \item The System from now on generates documents using the new filled in document template.
            \item The System stores the provided document template.
        \end{itemize}
        
    \item \textbf{Main scenario:} 
    \begin{enumerate}
       \item The Customer organization indicates it wants to download the current document template.
       \item The System looks up the correct document template and provides it to the Customer organization.
       \item The Customer organization changes the current document template.
       \item The Customer organization indicates it wants to upload the filled in template.
       \item The System downloads the updated document. It notifies an eDocs administrator that a new filled in document is waiting to be verified for correctness.
       \item An eDocs administrator verifies the correctness of the document template.
       \item The System stores it. The updated document replaces the old document template for that Customer organization and that document type.
    \end{enumerate}

    \item \textbf{Alternative scenarios:} 
    \begin{enumerate}
        \item [6a.] If the provided document template is not correctly filled in, the administrator notifies the Customer organization. Restart from step 3.
    \end{enumerate}
    
    \item \textbf{Remarks:}
        \begin{itemize}
            \item If the Customer organization has not filled in a document template before, the document template presented to the Customer organization is the document template that the eDocs administrator provided for the Customer Organization to use.
        \end{itemize}
\end{itemize}
\section{Non-functional requirements}\label{sec:non-functional}
In this section, we model the non-functional requirements for the system in the
form of \emph{quality attribute scenarios}. We provide for each type
(availability, performance and modifiability) one requirement.

\subsection{Availability}
\subsubsection{\emph{Av1}: Name of the quality attribute scenario}
Shortly describe the context of the scenario.

\begin{itemize}
    \item \textbf{Source:} source
    \item \textbf{Stimulus:}
        \begin{itemize}
            \item Description of a first stimulus.
            \item Description of a second stimulus.
        \end{itemize}

    \item \textbf{Artifact:} the stimulated artifact
    \item \textbf{Environment:} the condition under which the stimulus occurs
    \item \textbf{Response:}
        \begin{itemize}
            \item Describe how the system should respond to the stimulus.
        \end{itemize}

    \item \textbf{Response measure:}
        \begin{itemize}
            \item Describe how the satisfaction of a response is measured.
        \end{itemize}
\end{itemize}

\subsection{Performance}
\subsubsection{\emph{P1}: Name of the quality attribute scenario}
Shortly describe the context of the scenario.

\begin{itemize}
    \item \textbf{Source:} source
    \item \textbf{Stimulus:}
        \begin{itemize}
            \item Description of a first stimulus.
            \item Description of a second stimulus.
        \end{itemize}

    \item \textbf{Artifact:} the stimulated artifact
    \item \textbf{Environment:} the condition under which the stimulus occurs
    \item \textbf{Response:}
        \begin{itemize}
            \item Describe how the system should respond to the stimulus.
        \end{itemize}

    \item \textbf{Response measure:}
        \begin{itemize}
            \item Describe how the satisfaction of a response is measured.
        \end{itemize}
\end{itemize}

\subsection{Modifiability}
\subsubsection{\emph{M1}: Name of the quality attribute scenario}
Shortly describe the context of the scenario.

\begin{itemize}
    \item \textbf{Source:} source
    \item \textbf{Stimulus:}
        \begin{itemize}
            \item Description of a first stimulus.
            \item Description of a second stimulus.
        \end{itemize}

    \item \textbf{Artifact:} the stimulated artifact
    \item \textbf{Environment:} the condition under which the stimulus occurs
    \item \textbf{Response:}
        \begin{itemize}
            \item Describe how the system should respond to the stimulus.
        \end{itemize}

    \item \textbf{Response measure:}
        \begin{itemize}
            \item Describe how the satisfaction of a response is measured.
        \end{itemize}
\end{itemize}

\end{document}
